\section*{Vector Spaces: source notes}
\begin{enumerate}
\item In the proof of scalar identity for matrices, the final matrix contains sa, sb, sc, sd. Those entries must be a, b, c, d: the scalar is 1, not a free s.
\item In the closure-under-addition display for upper-triangular matrices, an equals sign is missing between the second addend and the resulting sum matrix.
\item For the operation r·(x,y)=(rx,0), the preceding distributivity counterexample does not work: both sides equal (0,0). The separately stated scalar-identity counterexample is valid and the answer “not a vector space” remains correct.
\item The requested arrowed-operation restatement is inconsistent in axioms (6)–(8): scalar multiplication still uses plain · in (6)–(7), and vector addition uses plain + in (8). These should use the newly defined arrowed scalar and vector operations; scalar addition r+s remains plain.
\item The span-membership row reduction omits a step. The two labelled row operations produce third row (0, 3/2, 7/2), not (0,0,3). Adding the resulting second row to the third then gives (0,0,3), so the “no solution” conclusion is unchanged.
\item In the three-dimensional superhero answer, the three device vectors span all of R³, not a plane. Their column matrix has determinant 18. The “no place to hide” conclusion and the displayed elimination are correct.
\item The span-extension proof should assume v belongs to span(S), not necessarily S. Its displayed linear-combination representation is exactly the weaker hypothesis needed by the proof.
\item The intersection proof changes its vector names: it assumes w and s belong to the intersection, then forms r v+s w. Use v and w consistently as vectors, with r and s as scalars.
\item The union-of-subspaces counterexample names ambient space R³ but displays two-entry coordinate vectors. Either use R² as the ambient space or append a zero third coordinate to each vector.
\item The condition given for span(S∪T)=span(S)∪span(T) is too strong. Equality holds exactly when one span contains the other, not only when one original set contains the other. For S=\{(1,0)\} and T=\{(2,0)\}, neither set contains the other but both spans are the same line. The following noncontainment argument confuses membership in a set with membership in its span.
\end{enumerate}
\section*{Linear Independence: source notes}
\begin{enumerate}
\item In the \{1+x,1-x\} example, R2 - R1 yields -2c2 = 0, not +2c2 = 0. The conclusion of independence is unchanged.
\item The general echelon-matrix statement needs the qualification nonzero rows. An echelon matrix with a zero row supplies an immediate counterexample. The displayed example itself has no zero row.
\item When isolating v from 0 = c1 s1 + ... + cn sn + c(n+1) v, each displayed coefficient of si needs a minus sign. The span-membership conclusion remains valid.
\item The first equation in the spanning-set reduction contains a plus sign in the empty c4 column, printing two successive plus signs. The intended first equation is c1 + c3 + 3c5 = 0.
\item In the polynomial exercise answer, the stated row operations give the last pivot -152/13, not -128/13. The independence conclusion remains correct.
\item In the three-row-vector exercise, the operation producing the zero third row must be R3 - (4/7)R2, not R3 + (4/7)R2. The final dependence vector is correct.
\item The union example defines S and T but its last sentence calls them S1 and S2. These are the same displayed sets; their union is dependent.
\item The question labels its proposed direction as if, but the proposed implication from an independent union to trivial intersection is the only-if direction. This note corrects the direction label only; it supplies no answer.
\item The supplied answer reverses the names if and only if: trivial intersection implies an independent union is the if direction; the reverse is only if. The corrected criterion and proof substance are unchanged.
\item The two-by-two determinant argument needs unique solution, not solution; every homogeneous system has a zero solution. Its combined fraction also needs denominator a, not d.
\item The three-by-three elimination step needs the negative multiplier -((ah-bg)/a) on row 2 to clear row 3. The displayed final pivot is the result of this subtraction.
\item In the a != 0, ae-bd = 0 case, the proof describes singularity while calling it nonsingularity. After the row swap the system is nonsingular exactly when both ah-bg and af-cd are nonzero, equivalently when the determinant is nonzero. Either factor zero, and the subsequent equal-zero calculation, establish singularity.
\item For the two vectors in R3, dependence requires ae-bd = ah-gb = dh-ge = 0. Equality of these three minors without the final zero condition is insufficient.
\end{enumerate}
\section*{Basis and Dimension: source notes}
\begin{enumerate}
\item The first polynomial-basis answer omits the contributions of c1 to the x and constant coefficients. Its proposed c2 and c3 therefore do not represent a general quadratic when a2 is nonzero. The stated set is still a basis. Correct coefficients (c1,c2,c3) are (a2; a0/2 + a1/4 - a2/4; -a0/2 + a1/4 + 3*a2/4).
\item The matrix-basis proof obtains c3=0 from the lower-right entries, not the lower-left entries as stated at line 918. The displayed third basis matrix has its sole nonzero entry at the lower right. The conclusion remains correct.
\item Choose s2 in S outside span(B1), and repeat that choice within S. Such an element exists whenever span(B1) is a proper subspace of span(S). B1 is initially a basis of its own span, not necessarily the full space.
\item Read the intermediate set as \{-1+x, 2x\textasciicircum{}2+x\textasciicircum{}3\}. The printed x\textasciicircum{}2+x\textasciicircum{}3 fails the exercise condition a2-2a3=0; the final displayed basis and dimension are correct.
\item The two displayed sets describe the same diagonal-matrix space; an equality sign is missing at the start of the second line. Preserve the original display and identify this separately.
\item Use a nontrivial relation among the concatenated bases. Independence within each basis implies that the partial combination from either basis is nonzero. These equal, opposite-signed partial combinations provide a nonzero vector in U intersect W; no individual basis vector need lie in the other space.
\item Read V=W in the first sentence of this part as U=W. The remainder of the argument already uses U and W.
\item The right-hand column in the first displayed equation should be (1,3), as in the question and subsequent system. The conclusion that it is outside the column space remains correct.
\item At the second arrow, read the row-four operation as -2 times row one plus row four. The displayed echelon matrix then follows. The later basis ending in -5x\textasciicircum{}2 is valid: scaling the nonzero third basis vector does not change its span.
\item Read “at least two” here as “at most two”. The fourth and fifth columns are independent, so the column rank is exactly two, as stated.
\item For unequal row counts, first append zero rows to the shorter matrix, or compare only the nonzero reduced rows. The exercise claims equality of nonzero rows, not equality of differently sized matrices. A proof can stack A over B and B over A: the two stacked matrices are row equivalent, and reducing either dependent extra block leaves the same nonzero rows.
\item In part (b), read the right-hand entries as d1 through dm. There is one right-hand entry for each of the m equations.
\item For the missing direction, if every right-hand vector in R\textasciicircum{}m is attainable, the columns span R\textasciicircum{}m. The column rank is then m, and equality of row and column ranks gives full row rank. The second printed paragraph is the contrapositive of the first direction.
\item In part (e), W1 is a subspace: shifting its parameter still describes the entire x-axis. W2 is not a subspace because it does not contain zero. The answer “No” remains correct.
\item Read “greater than” here as “at least”. Dimensions eight, nine and ten for the sum, and eight, seven and six for the intersection, are all possible. The printed final list is correct.
\item Read “every subset” as “every subspace”. Closure under addition is available because W is a subspace; the identity W+W=W is correct.
\end{enumerate}
\section*{Fields: source notes}
\begin{enumerate}
\item The displayed axiom list does not explicitly require the additive and multiplicative identities to be distinct. Its literal conditions also hold in the singleton algebra with 0=1; the usual field definition excludes this case. Record the missing explicit condition separately, without changing the pinned source.
\item The supplied binary-field answer calls addition closure and commutativity the first half of condition (2), but those appear in condition (1). The operations and the finite-field conclusion are correct.
\end{enumerate}
